%% LyX 2.4.4 created this file.  For more info, see https://www.lyx.org/.
%% Do not edit unless you really know what you are doing.
\documentclass[english,t]{beamer}
\usepackage{mathptmx}
\usepackage[T1]{fontenc}
\usepackage[latin9]{inputenc}
\usepackage{amssymb}
\usepackage{cancel}

\makeatletter
%%%%%%%%%%%%%%%%%%%%%%%%%%%%%% Textclass specific LaTeX commands.
% this default might be overridden by plain title style
\newcommand\makebeamertitle{\frame{\maketitle}}%
% (ERT) argument for the TOC
\AtBeginDocument{%
  \let\origtableofcontents=\tableofcontents
  \def\tableofcontents{\@ifnextchar[{\origtableofcontents}{\gobbletableofcontents}}
  \def\gobbletableofcontents#1{\origtableofcontents}
}

%%%%%%%%%%%%%%%%%%%%%%%%%%%%%% User specified LaTeX commands.
\usetheme{Warsaw}
% or ...

\setbeamercovered{transparent}
% or whatever (possibly just delete it)
\setbeamercovered{invisible} %makes overlays invisible


%gets rid of navigation symbols
\setbeamertemplate{navigation symbols}{}

\usepackage{multicol}

\usepackage{tikz}
\usetikzlibrary{calc}
\usetikzlibrary{plotmarks}
\usepackage{pgfplots}
\pgfplotsset{compat=newest}

\usepackage{calc}
\usepackage[nomessages]{fp}

\usepackage{advdate}    % Advancing/saving dates
\usepackage{datetime}   % Dates formatting
\usepackage{datenumber} % Counters for dates
\usepackage[super]{nth}

%\usepackage{pgfpages}
%\pgfpagesuselayout{4 on 1}[a4paper, landscape, border shrink=7mm]
%\pgfpageslogicalpageoptions{1}{border code=\pgfusepath{stroke}}
%\pgfpageslogicalpageoptions{2}{border code=\pgfusepath{stroke}}
%\pgfpageslogicalpageoptions{3}{border code=\pgfusepath{stroke}}
%\pgfpageslogicalpageoptions{4}{border code=\pgfusepath{stroke}}
%%%Don't forget to add 'handout'

\makeatother

\usepackage{babel}
\begin{document}
\newcommand{\xmax}{2000}
\newcommand{\xmin}{0}
\newcommand{\xstep}{0} %must be xmin+n
\newcommand{\ymax}{500}
\newcommand{\ymin}{0}
\newcommand{\ystep}{0} %must be ymin+n

\newcounter{countEx} \setcounter{countEx}{0}
\newcounter{countProb} \setcounter{countProb}{0}
\resetcounteronoverlays{countEx} \resetcounteronoverlays{countProb}

\newcommand{\ex}[3]{
	\addtocounter{countEx}{1}
	\only<#1-#2>{\begin{exampleblock} {Example \chNum.\secNum.\arabic{countEx}.} #3	\end{exampleblock}}
}

\newcommand{\prob}[4]{
	\addtocounter{countProb}{1}
	\only<#1-#2>{\begin{exampleblock} {Problem  \chNum.\secNum.\arabic{countProb}.\,#3} #4	\end{exampleblock}}
}

\newcommand{\yline}[4]{
	(#2 - #4)/(#1 - #3)*(x - #1) + #2
}

\beamerdefaultoverlayspecification{<+->}

%%%%Chapter and Section numbers%%%%%
\newcommand{\chNum}{5}
\newcommand{\secNum}{1}
\newcommand{\secTitle}{Simple and Compound Interest}

\title[Section \chNum.\secNum: \secTitle]{Section \chNum.\secNum: \secTitle}

\subtitle{Finite Mathematics~~MTH 132~~Spring 2014}

\author{Dr.\,James Ashe}

\titlegraphic{\includegraphics[height=1.5cm]{/home/jim/JamesAshe/JCSU/images/jcsu_bull}}

\institute{Johnson C. Smith University}

\date{{}}

\makebeamertitle 
\begin{frame}{Mathematics of Finance}

\begin{columns}


\column{10cm}
\begin{itemize}
\item [Chapter 5:]\emph{ Mathematics of Finance}

\begin{itemize}
\item [5.1] Simple and Compound Interest
\item [5.2] Future Value of an Annuity
\item [5.3] Present Value of an Annuity; Amortization
\end{itemize}
\end{itemize}
\end{columns}

\begin{columns}


\column{3.5cm}
\begin{exampleblock}<2->{Example 5.1}
 If you borrow \$7200, at an annual interest rate of 6.2\% compounded
daily, how much is owed after 15 months?
\end{exampleblock}

\column{3.5cm}
\begin{exampleblock}<3->{Example 5.2}
 You need to save \$250,000 for retirement in 20 years. What size
should the monthly payments into a fund which pays 7\% compounded
monthly be?
\end{exampleblock}

\column{3.5cm}
\begin{exampleblock}<4->{Example 5.3}
You have \$8,500 in credit card debt at 25\% compounded monthly.
If you want to pay off the debt on 5 years, how large should the monthly
payments be?
\end{exampleblock}
\end{columns}
\end{frame}

\begin{frame}{Simple Interest}

\begin{columns}


\column{6cm}

\ex{2}{}{John borrows \$2,000 at 6.3\% for 39 weeks. How much interest did he accrue?}

\column{5cm}
\begin{block}{Simple Intererst}
$I=Prt$

\begin{itemize}
\item []$P$ is the \emph{principal}
\item []$r$ is the annual interest \emph{rate} 
\item []$t$ is the \emph{time} in years
\end{itemize}
\end{block}
\end{columns}
\end{frame}

\begin{frame}{Simple Interest}

\begin{columns}


\column{6cm}

\ex{1}{}{John borrows \$2,000 at 6.3\% for 39 weeks. How much does he owe?}

\column{5cm}
\begin{block}<2->{Future/Maturity Value for Simple Interest}
$A=P(1+rt)$

\begin{itemize}
\item []$A$ is the future\emph{ amount.}
\end{itemize}
\end{block}
\end{columns}
\end{frame}

\begin{frame}{Simple Interest}

\begin{columns}


\column{6cm}

\ex{1}{}{A CD pays 3.75\% after 5 years. If John wants to get \$5,000 from the CD how much should he invest?}

\column{5cm}
\begin{block}<2->{Present Value for Simple Interest}
$P=\frac{A}{1+rt}$
\end{block}
\end{columns}
\end{frame}

\begin{frame}{Compound Interest}

\begin{columns}


\column{6cm}

\ex{1}{}{John borrows \$2,000 at 6.3\% \emph{compounded monthly} for 39 weeks. How much does he owe?}

\column{5cm}
\begin{block}<2->{Future/Maturity Value for Compund Interest}
$A=P(1+i)^{n}$

\begin{itemize}
\item []$i=\frac{r}{m}$ and $n=mt$
\item []$m$ is the number of \emph{compounding periods per year.}
\item []$n$ is the number of \emph{compounding periods}. 
\item []$i$ is the \emph{interest rate per period}.
\end{itemize}
\end{block}
\end{columns}
\end{frame}

\begin{frame}{Compound Interest}

\begin{columns}


\column{6cm}

\ex{1}{}{An account pays 3.75\% compounded semi-annually. If John wants to get \$5,000 in 5 years from the account how much should he invest?}

\column{5cm}
\begin{block}<2->{Present Value for Compound Interest}
$P=\frac{A}{(1+i)^{n}}$

\begin{itemize}
\item []$i=\frac{r}{m}$ and $n=mt$
\item []$m$ is the number of \emph{compounding periods per year.}
\item []$n$ is the number of \emph{compounding periods}. 
\item []$i$ is the \emph{interest rate per period}.
\end{itemize}
\end{block}
\end{columns}
\end{frame}

\begin{frame}{Determining $m$}

\begin{columns}


\column{10cm}
\begin{itemize}
\item Yearly\uncover<2->{$\Rightarrow m=1$}
\item Biannually/semi-annually\uncover<3->{$\Rightarrow m=2$}
\item Quarterly\uncover<4->{$\Rightarrow m=4$}
\item Monthly\uncover<5->{$\Rightarrow m=12$}
\item Weekly\uncover<5->{$\Rightarrow m=52$}
\item Daily\uncover<6->{$\Rightarrow m=365$ for \emph{exact }or $360$
for \emph{ordinary interest}.}
\item Continuously\uncover<6->{$\Rightarrow m=\infty$}
\end{itemize}
\end{columns}
\end{frame}

\begin{frame}{Continuous Compounding}

\begin{columns}


\column{6cm}

\ex{1}{}{John borrows \$2,000 at 6.3\% \emph{compounded continuously} for 39 weeks. How much does he owe?}

\column{5cm}
\begin{block}{The number \$e\$}
\begin{eqnarray*}
e & = & {\displaystyle \lim_{n\rightarrow\infty}\left(1+\frac{1}{n}\right)^{n}}\\
 & \approx & 2.718281
\end{eqnarray*}
\end{block}

\begin{block}<2->{Future/Maturity Value for Continuously Compound Interest}
$A=Pe^{rt}$
\end{block}
\end{columns}
\end{frame}

\begin{frame}{Continuous Compounding}

\begin{columns}


\column{6cm}

\ex{1}{}{John owes \$4,000 at 5\% \emph{compounded continuously} after 10 years. How much did he borrow?}

\column{5cm}

\begin{block}<2->{Present Value for Continuously Compound Interest}
$P=\frac{A}{e^{rt}}=Ae^{-rt}$
\end{block}
\end{columns}
\end{frame}

\begin{frame}{Effective Rate}

\small{If John borrows $\$1$ compounded daily at a rate of $5\%$
for one year he will owe
\begin{itemize}
\item <2->$1\cdot(1+\frac{.05}{365})^{1\cdot365}\approx1.0513$
\item <3->So borrowing $\$1$ actually cost John $\$0.0513$ NOT $\$0.05$
as advertised.
\end{itemize}
\begin{block}<4->{}
 The \emph{effective rate }corresponding to a stated rate of interest
$r$ compounded $m$ times per year is 
\[
r_{E}=\left(1+\frac{r}{m}\right)^{m}-1
\]
\end{block}

\begin{block}<5->{}
The \emph{effective rate }corresponding to a stated rate of interest
$r$ compounded continuously is 
\[
r_{E}=e^{r}-1
\]
\end{block}
\begin{itemize}
\item <6->In the ``real world'' effective rate is usually called \emph{APR
}or \emph{APY.}
\end{itemize}
}
\end{frame}

\end{document}
